
\begin{abstract}
High-fidelity three-dimensional (3-D) computational fluid dynamics (CFD) data are expensive to generate, limiting the training of deep-learning models for transonic finite-wing surface-pressure prediction. Lower-dimensional aerodynamic data provide transferable sectional information but cannot fully represent finite-span effects, particularly near the wing root and tip. To address this limitation, this study develops a hierarchical aerodynamic knowledge-transfer framework that separates sectional pressure responses from localized three-dimensional corrections. A baseline pressure model is trained on infinite-span swept-wing data, while root- and tip-correction models are developed using pressure differences between paired finite- and infinite-span CFD solutions. These knowledge models are transferred to the target wing to construct a three-dimensional pressure prior, which is subsequently refined through residual learning using limited target-domain CFD samples. With only 10 target-wing training samples, the proposed method reduces the surface-pressure root-mean-square error by 65.74\% relative to direct three-dimensional learning and achieves lower prediction error than the direct-learning model trained with 50 samples. Ablation studies further demonstrate the complementary contributions of root and tip corrections. In a downstream aerodynamic-coefficient regression task using 20 force-labeled wings, incorporating pressure distributions predicted by a separately trained model reduces the mean absolute errors in lift and drag coefficients by 70.1\% and 54.8\%, respectively, compared with direct regression. These results demonstrate that combining transferable sectional knowledge with localized finite-span corrections provides an effective aerodynamic prior, substantially improving the data efficiency of three-dimensional wing pressure prediction.
\end{abstract}
